\documentclass[aspectratio=169, 10pt]{beamer}

% --- Idioma e Codificação ---
\usepackage[utf8]{inputenc}
\usepackage[T1]{fontenc}
\usepackage[brazil]{babel}
\usepackage{lmodern}

% --- Pacotes Gráficos e Diagramação ---
\usepackage{graphicx}
\usepackage{booktabs}
\usepackage{tikz}
\usetikzlibrary{shapes,arrows,positioning,calc}
\usepackage{listings}
\usepackage{tcolorbox}
\tcbuselibrary{skins,breakable}
\usepackage{qrcode}

% --- Cores Institucionais do IFF e do CONEPE ---
\definecolor{iffgreen}{RGB}{0, 102, 51}         % Verde IFF
\definecolor{iffgreenlight}{RGB}{235, 247, 238}
\definecolor{iffdark}{RGB}{15, 60, 30}
\definecolor{iffred}{RGB}{204, 0, 0}            % Vermelho IFF
\definecolor{codebg}{RGB}{245, 247, 248}
\definecolor{conepeteal}{RGB}{43, 114, 131}     % Tom azul-petróleo do CONEPE 2026

% --- Tema do Beamer ---
\usetheme{Madrid}
\useinnertheme{rectangles}

% --- DIRETRIZ 1: SEM RODAPÉ ---
\setbeamertemplate{navigation symbols}{}
\setbeamertemplate{footline}{} % Remove completamente o rodapé

% --- DIRETRIZ 2: CABEÇALHO COM O BANNER DO CONEPE 2026 ---
\setbeamertemplate{frametitle}{%
  \vspace{-1pt}%
  \begin{beamercolorbox}[wd=\paperwidth,ht=1.15cm,dp=0ex]{empty}%
    \includegraphics[width=\paperwidth,height=1.15cm]{cabecalho_conepe2026.jpg}%
  \end{beamercolorbox}%
  \vspace{0.15cm}%
  \begin{beamercolorbox}[wd=\paperwidth,leftskip=0.8cm,rightskip=0.8cm]{frametitle}%
    \usebeamerfont{frametitle}\color{iffdark}\textbf{\insertframetitle}%
    \ifx\insertframesubtitle\@empty%
    \else%
      \vskip0.1em{\usebeamerfont{framesubtitle}\color{gray}\insertframesubtitle}%
    \fi%
  \end{beamercolorbox}%
}

% --- Customização de Cores dos Blocos ---
\setbeamercolor{structure}{fg=iffgreen}
\setbeamercolor{block title}{bg=iffgreen,fg=white}
\setbeamercolor{block body}{bg=iffgreenlight,fg=black}
\setbeamercolor{block title alerted}{bg=iffred,fg=white}
\setbeamercolor{block body alerted}{bg=iffred!10,fg=black}

% --- Configurações de Código TeX ---
\lstset{
  basicstyle=\ttfamily\scriptsize,
  keywordstyle=\color{iffgreen}\bfseries,
  commentstyle=\color{gray}\itshape,
  stringstyle=\color{iffred},
  backgroundcolor=\color{codebg},
  frame=single,
  rulecolor=\color{gray!30},
  breaklines=true,
  showstringspaces=false,
  tabsize=2
}

% --- Metadados do Trabalho ---
\hypersetup{
  pdftitle={Classe LaTeX: Otimizacao da Formatacao de Trabalhos Academicos - CONEPE 2026},
  pdfauthor={P. H. R. De Andrade, A. C. Soja, M. L. L. Dantas, A. M. O. Figueiredo}
}

\title[\texttt{ifftese.cls} --- CONEPE 2026]{Classe \LaTeX: Otimização da Formatação de Trabalhos Acadêmicos}
\subtitle{Automação Normativa da ABNT e Abstração Estrutural com \texttt{ifftese.cls} e \texttt{macros.sty}}

\author[P. H. R. De Andrade \textit{et al.}]{
  \textbf{P. H. R. De Andrade}\textsuperscript{1*}, A. C. Soja\textsuperscript{1}, M. L. L. Dantas\textsuperscript{2}, A. M. O. Figueiredo\textsuperscript{1}
}

\institute[IFF / PUC-Chile]{
  \textsuperscript{1}Instituto Federal Fluminense (IFF) --- \textit{Campus} Bom Jesus do Itabapoana\\
  \textsuperscript{2}Pontifícia Universidade Católica de Chile (PUC-Chile)\\[1ex]
  \texttt{*pedroiff0@gmail.com}\\[1.5ex]
  \footnotesize \textbf{Apoio Institucional e Financeiro:}\\
  \textbf{CNPq} (Conselho Nacional de Desenvolvimento Científico e Tecnológico) --- Projeto nº 129985/2025-2\\
  \textbf{Instituto Federal Fluminense} --- \textit{Campus} Bom Jesus do Itabapoana
}

\date{CONEPE 2026 --- 22 a 24 de Setembro de 2026}

% -----------------------------------------------------------------------------
% DOCUMENTO
% -----------------------------------------------------------------------------
\begin{document}

% --- Slide 1: Capa Oficial ---
{
\setbeamertemplate{frametitle}{}
\begin{frame}[plain]
  \begin{beamercolorbox}[wd=\paperwidth,ht=1.3cm,dp=0ex]{empty}
    \includegraphics[width=\paperwidth,height=1.3cm]{cabecalho_conepe2026.jpg}
  \end{beamercolorbox}
  \vspace{-0.2cm}
  \titlepage
\end{frame}
}

% --- Slide 2: Storytelling - O Pesadelo no Word ---
\begin{frame}{O Pesadelo da Madrugada: A Odisseia da Formatação no Word}
  \begin{columns}[T]
    \begin{column}{0.55\textwidth}
      \textbf{Cenário:} Véspera da entrega final, 23h59.
      \begin{itemize}
        \item O orientador solicita a inclusão de um parágrafo de 3 linhas na Introdução.
        \item Você insere o texto e pressiona \textit{Enter}...
      \end{itemize}
      \vspace{0.2cm}
      \begin{alertblock}{O Efeito Dominó Catastrófico}
        \begin{itemize}
          \item A \textbf{Figura 4} na pág. 8 salta misteriosamente para a pág. 11, deixando um enorme vazio branco.
          \item A legenda da figura se desprende e fica órfã no topo da folha seguinte.
          \item Ocorre o temido: \texttt{"Erro! Indicador não definido."}
          \item As referências numéricas ficam fora de ordem: \texttt{[14, 2, 8, 1]}.
        \end{itemize}
      \end{alertblock}
    \end{column}

    \begin{column}{0.42\textwidth}
      \begin{tcolorbox}[colback=red!5,colframe=iffred,title=\textbf{Sintomas do WYSIWYG}]
        \small
        \begin{itemize}
          \item Quebras de seção corrompem a paginação pré-textual.
          \item Margens que mudam sozinhas.
          \item Sumário desalinha os números de página dos pontilhados (\texttt{.....}).
          \item \textbf{Custo real:} 70\% do tempo gasto brigando contra a ferramenta, não produzindo ciência.
        \end{itemize}
      \end{tcolorbox}
    \end{column}
  \end{columns}
\end{frame}

% --- Slide 3: O Paradoxo do Pesquisador ---
\begin{frame}{O Paradoxo do Pesquisador Acadêmico}
  \begin{columns}[c]
    \begin{column}{0.5\textwidth}
      \begin{tcolorbox}[colback=codebg,colframe=conepeteal,title=\textbf{O Objetivo Central}]
        \begin{itemize}
          \item Formulação de hipóteses científicas;
          \item Coleta e análise rigorosa de dados;
          \item Redação clara e impacto na sociedade;
          \item Contribuição para o avanço do conhecimento.
        \end{itemize}
      \end{tcolorbox}
    \end{column}

    \begin{column}{0.5\textwidth}
      \begin{tcolorbox}[colback=orange!10,colframe=orange!80!black,title=\textbf{A Realidade Operacional}]
        \begin{itemize}
          \item Ajustar espaçamento de 1,5 cm vs 1,0 cm;
          \item Medir margem de 3 cm e 2 cm na régua visual;
          \item Contar páginas sem imprimir na folha de rosto;
          \item Repadronizar títulos após o Word perder estilos.
        \end{itemize}
      \end{tcolorbox}
    \end{column}
  \end{columns}
  \vspace{0.4cm}
  \centering \textbf{Pergunta norteadora:} \textit{Como blindar o pesquisador contra essas falhas sem exigir dele conhecimento avançado em programação?}
\end{frame}

% --- Slide 4: Normas Regentes I (NBR 14724) ---
\begin{frame}{As Normas Regentes I: ABNT NBR 14724 (Trabalhos Acadêmicos)}
  \begin{columns}[T]
    \begin{column}{0.48\textwidth}
      \begin{block}{\textbf{Geometria e Margens Estritas}}
        \begin{itemize}
          \item \textbf{Anverso:} Esquerda e Superior: \textbf{3 cm}; Direita e Inferior: \textbf{2 cm}.
          \item \textbf{Verso:} Margens espelhadas.
          \item \textbf{Fonte:} Tamanho 12 no corpo; tamanho menor (10) para citações longas, notas e fontes.
        \end{itemize}
      \end{block}
    \end{column}

    \begin{column}{0.48\textwidth}
      \begin{alertblock}{\textbf{A Pegadinha da Paginação}}
        \begin{itemize}
          \item As folhas pré-textuais são \textbf{contadas} a partir da Folha de Rosto;
          \item Porém, a numeração só pode ser \textbf{impressa a partir da Introdução};
          \item No Word: exige quebras de seção delicadas e desvinculação de cabeçalho;
          \item Na \texttt{ifftese.cls}: gerido automaticamente na transição textual.
        \end{itemize}
      \end{alertblock}
    \end{column}
  \end{columns}
  \vspace{0.2cm}
  \begin{tcolorbox}[colback=codebg,colframe=gray!50]
    \centering \footnotesize \textbf{Citações Diretas Longas (> 3 linhas):} Recuo obrigatório de 4 cm da margem esquerda, fonte reduzida e espaçamento simples.
  \end{tcolorbox}
\end{frame}

% --- Slide 5: Normas Regentes II (NBR 6023 e 6027) ---
\begin{frame}{As Normas Regentes II: NBR 6023 (Referências) e NBR 6027 (Sumário)}
  \begin{columns}[T]
    \begin{column}{0.48\textwidth}
      \begin{block}{\textbf{ABNT NBR 6023 (Referências)}}
        \begin{itemize}
          \item \textbf{Alinhamento:} À margem esquerda (nunca justificado!).
          \item \textbf{Espaçamento:} Entrelinhas simples; separadas entre si por 1 espaço simples em branco.
          \item \textbf{Destaque Tipográfico:} Uniforme em todo o documento (negrito ou itálico) estritamente no título da obra.
          \item \textbf{Sistema:} Autor-data ou numérico consistente.
        \end{itemize}
      \end{block}
    \end{column}

    \begin{column}{0.48\textwidth}
      \begin{block}{\textbf{ABNT NBR 6027 (Sumário)}}
        \begin{itemize}
          \item \textbf{Identidade Tipográfica:} As seções devem ter no sumário a exata grafia do texto:
            \begin{itemize}
              \item 1 SEÇÃO PRIMÁRIA (CAIXA ALTA NEGRITO)
              \item 1.1 Seção Secundária (Caixa Baixa Normal)
            \end{itemize}
          \item \textbf{Alinhamento:} Números de página ligados por linhas contínuas pontilhadas (\texttt{.....}).
        \end{itemize}
      \end{block}
    \end{column}
  \end{columns}
\end{frame}

% --- Slide 6: Padrão IBGE (Tabelas vs. Quadros) ---
\begin{frame}{Diferenciação Normativa: Tabelas (IBGE) vs. Quadros (ABNT)}
  \begin{columns}[T]
    \begin{column}{0.48\textwidth}
      \begin{tcolorbox}[colback=iffgreenlight,colframe=iffgreen,title=\textbf{Tabelas (Normas do IBGE)}]
        \begin{itemize}
          \item Finalidade: Predominância de \textbf{dados numéricos}.
          \item Estrutura: \textbf{Laterais abertas} (sem traços verticais externos de fechamento).
          \item Delimitação: Linha horizontal no topo, no cabeçalho e na base.
          \item Pré-texto: Conduz à \textbf{Lista de Tabelas}.
        \end{itemize}
      \end{tcolorbox}
    \end{column}

    \begin{column}{0.48\textwidth}
      \begin{tcolorbox}[colback=iffred!10,colframe=iffred,title=\textbf{Quadros (Norma ABNT)}]
        \begin{itemize}
          \item Finalidade: Dados \textbf{qualitativos, esquemáticos ou textuais}.
          \item Estrutura: \textbf{Totalmente fechado} nos quatro lados, com divisórias de células.
          \item Pré-texto: Exige obrigatoriamente a criação da \textbf{Lista de Quadros}.
        \end{itemize}
      \end{tcolorbox}
    \end{column}
  \end{columns}
  \vspace{0.2cm}
  \begin{alertblock}{O Erro Típico na Academia}
    Estudantes tratam quadros e tabelas como sinônimos. A \texttt{ifftese.cls} automatiza a separação através das macros dedicadas \texttt{\textbackslash inserirtabela} e \texttt{\textbackslash inserirquadro}.
  \end{alertblock}
\end{frame}

% --- Slide 7: Princípio dos Paradigmas ---
\begin{frame}{O Princípio dos Paradigmas: WYSIWYG vs. WYSIWYM}
  \begin{columns}[T]
    \begin{column}{0.48\textwidth}
      \begin{tcolorbox}[colback=red!5,colframe=iffred,title=\textbf{WYSIWYG: Microsoft Word}]
        \textit{"What You See Is What You Get"}
        \begin{itemize}
          \item Mistura conteúdo, apresentação e geometria simultaneamente.
          \item Algoritmo de renderização local ("guloso"): a edição de uma linha propaga efeitos colaterais por todo o arquivo.
          \item Instável para teses longas (> 50 páginas) com muitas ilustrações e equações.
        \end{itemize}
      \end{tcolorbox}
    \end{column}

    \begin{column}{0.48\textwidth}
      \begin{tcolorbox}[colback=iffgreenlight,colframe=iffgreen,title=\textbf{WYSIWYM: O Sistema \LaTeX}]
        \textit{"What You See Is What You Mean"}
        \begin{itemize}
          \item \textbf{Donald Knuth (1977)} e \textbf{Leslie Lamport (1984)}.
          \item Separação ontológica: o autor declara a semântica (\texttt{\textbackslash section}, \texttt{\textbackslash caption}); o compilador calcula o layout ótimo.
          \item Algoritmo de Knuth-Plass: microtipografia globalmente calculada.
        \end{itemize}
      \end{tcolorbox}
    \end{column}
  \end{columns}
\end{frame}

% --- Slide 8: O Dilema do LaTeX Puro ---
\begin{frame}{Por Que o \LaTeX\ é a Resposta (e por que ainda causa receio?)}
  \begin{columns}[c]
    \begin{column}{0.5\textwidth}
      \textbf{As Vantagens Inquestionáveis:}
      \begin{itemize}
        \item Padrão ouro nos repositórios arXiv e NASA/ADS;
        \item Tipografia matemática perfeita e estável;
        \item Referências cruzadas confiáveis via arquivos \texttt{.aux};
        \item Documento imune a quebras bruscas de layout.
      \end{itemize}
    \end{column}

    \begin{column}{0.5\textwidth}
      \begin{alertblock}{\textbf{A Barreira de Entrada}}
        \begin{itemize}
          \item Curva de aprendizado íngreme;
          \item Sintaxe densa com dezenas de pacotes gráficos;
          \item Erros de compilação assustadores para o discente iniciante;
          \item Dificuldade em adequar classes genéricas às normas da ABNT.
        \end{itemize}
      \end{alertblock}
    \end{column}
  \end{columns}
  \vspace{0.4cm}
  \centering \textbf{A Solução Necessária:} Uma camada de abstração projetada sob medida para a nossa instituição.
\end{frame}

% --- Slide 9: A Solução IFF ---
\begin{frame}{A Solução Proposta: \texttt{ifftese.cls} e \texttt{macros.sty}}
  \begin{block}{\textbf{Objetivo do Projeto}}
    Desenvolver uma classe tipográfica oficial e um pacote de extensão encapsulado para atender a comunidade do Instituto Federal Fluminense, assegurando conformidade estrita à ABNT e preservando a identidade visual local.
  \end{block}

  \vspace{0.3cm}

  \begin{columns}[T]
    \begin{column}{0.48\textwidth}
      \begin{tcolorbox}[colback=white,colframe=iffgreen,title=\textbf{Classe \texttt{ifftese.cls}}]
        \small
        \begin{itemize}
          \item Gerenciamento de margens e geometria da folha;
          \item Geração automática de capas e folhas de aprovação;
          \item Cabeçalhos institucionais e metadados formais;
          \item Base sólida construída sobre o \texttt{abntex2}.
        \end{itemize}
      \end{tcolorbox}
    \end{column}

    \begin{column}{0.48\textwidth}
      \begin{tcolorbox}[colback=white,colframe=conepeteal,title=\textbf{Extensão \texttt{macros.sty}}]
        \small
        \begin{itemize}
          \item Abstração de comandos complexos em linha única;
          \item Diferenciação nativa de tabelas e quadros;
          \item Blindagem contra falhas em listas dinâmicas;
          \item Manipulação segura de contadores pós-textuais.
        \end{itemize}
      \end{tcolorbox}
    \end{column}
  \end{columns}
\end{frame}

% --- Slide 10: Metodologia de Desenvolvimento ---
\begin{frame}{Metodologia de Desenvolvimento em Três Etapas}
  \centering
  \begin{tikzpicture}[node distance=2.4cm, auto, >=latex]
    \tikzstyle{block} = [rectangle, draw=iffgreen, fill=iffgreenlight, text width=3.8cm, text centered, rounded corners, minimum height=1.3cm, font=\small\bfseries]
    \tikzstyle{line} = [draw, -latex', thick, iffgreen]

    \node [block] (step1) {1. Mapeamento Normativo\\{\footnotesize NBR 14724, 6023, 6027 e IBGE}};
    \node [block, right of=step1, xshift=2.2cm] (step2) {2. Construção das Macros\\{\footnotesize Encapsulamento de ambientes}};
    \node [block, right of=step2, xshift=2.2cm] (step3) {3. Template Unificado\\{\footnotesize Modelo \texttt{main.tex} pronto para uso}};

    \path [line] (step1) -- (step2);
    \path [line] (step2) -- (step3);
  \end{tikzpicture}

  \vspace{0.6cm}
  \begin{columns}[T]
    \begin{column}{0.32\textwidth}
      \footnotesize \textbf{Ambientes de Teste:} TeX Live com motores \texttt{pdflatex} e \texttt{bibtex}.
    \end{column}
    \begin{column}{0.32\textwidth}
      \footnotesize \textbf{Acessibilidade:} Compatível com Overleaf, TeXPage e repositórios locais.
    \end{column}
    \begin{column}{0.32\textwidth}
      \footnotesize \textbf{Engenharia de Software:} Foco centrado no usuário final e experiência do autor.
    \end{column}
  \end{columns}
\end{frame}

% --- Slide 11: Resultados - Variáveis e Pré-Texto ---
\begin{frame}[fragile]{Resultados: Variáveis Semânticas e Automação Pré-Textual}
  \begin{columns}[T]
    \begin{column}{0.48\textwidth}
      \textbf{Switches de Configuração (Flags Booleanas):}
      \begin{itemize}
        \item \verb|\frenteVerso|: Alterna margens simples ou espelhadas.
        \item \verb|\capaiff|: Aplica leiaute institucional oficial.
        \item \verb|\corlink|: Controle de contraste para impressão ou tela.
        \item \verb|\sumarioEscada|: Formatação hierárquica automática.
      \end{itemize}
      \vspace{0.2cm}
      \textbf{Metadados Simples no Preâmbulo:}
      \begin{itemize}
        \item \verb|\autor|, \verb|\titulo|, \verb|\orientador|, \verb|\local|, \verb|\data|.
      \end{itemize}
    \end{column}

    \begin{column}{0.5\textwidth}
      \begin{tcolorbox}[colback=codebg,colframe=iffdark,title=\textbf{Geração com Comandos Únicos}]
        \small
        O usuário não formata a capa manualmente:
        \begin{lstlisting}[language=TeX]
\capa
\contracapa
\folhadeaprovacao
        \end{lstlisting}
        \vspace{0.2cm}
        \textbf{Garantias do Compilador:}
        \begin{itemize}
          \item Espaçamentos milimétricos exatos da NBR 14724;
          \item Contagem oculta de páginas ativada;
          \item Zero intervenção em arquivos auxiliares.
        \end{itemize}
      \end{tcolorbox}
    \end{column}
  \end{columns}
\end{frame}

% --- Slide 12: Resultados - Elementos Textuais: Inserção de Figuras ---
\begin{frame}[fragile]{Resultados: Elementos Textuais (\texttt{\textbackslash inserirfigura})}
  \begin{columns}[T]
    \begin{column}{0.48\textwidth}
      \begin{tcolorbox}[colback=red!5,colframe=iffred,title=\textbf{\LaTeX\ Convencional (8--10 linhas)}]
\begin{lstlisting}[language=TeX]
\begin{figure}[htbp]
  \centering
  \caption{Diagrama de Fases}
  \includegraphics[width=0.7\textwidth]{fases.png}
  \vspace{0.1cm}
  {\footnotesize Fonte: Autor (2026).}
  \label{fig:fases}
\end{figure}
\end{lstlisting}
        \scriptsize \textbf{Riscos:} Esquecimento de centralização, inversão de legenda/fonte, esquecimento do \texttt{\textbackslash label}.
      \end{tcolorbox}
    \end{column}

    \begin{column}{0.48\textwidth}
      \begin{tcolorbox}[colback=iffgreenlight,colframe=iffgreen,title=\textbf{Com o \texttt{macros.sty} (1 linha!)}]
\begin{lstlisting}[language=TeX]
\inserirfigura[0.7]{fases.png}
  {Diagrama de Fases}
  {Autor (2026)}{fig:fases}
\end{lstlisting}
        \vspace{0.4cm}
        \footnotesize \textbf{Garantias Automáticas:}
        \begin{itemize}
          \item Legenda rigorosamente no topo (norma ABNT);
          \item Fonte em tamanho menor no rodapé;
          \item Centralização horizontal padronizada;
          \item Geração dinâmica do rótulo e envio à lista.
        \end{itemize}
      \end{tcolorbox}
    \end{column}
  \end{columns}
\end{frame}

% --- Slide 13: Resultados - Automação IBGE e Pós-Texto ---
\begin{frame}{Resultados: Automação IBGE e Blindagem do Pós-Texto}
  \begin{columns}[T]
    \begin{column}{0.48\textwidth}
      \begin{block}{\textbf{Tabelas e Quadros Parametrizados}}
        \begin{itemize}
          \item \texttt{\textbackslash inserirtabela\{arq\}\{legenda\}\{fonte\}\{lbl\}}
            \begin{itemize}
              \item Gera laterais abertas conforme o IBGE;
              \item Alimenta a Lista de Tabelas.
            \end{itemize}
          \item \texttt{\textbackslash inserirquadro\{arq\}\{legenda\}\{fonte\}\{lbl\}}
            \begin{itemize}
              \item Gera moldura fechada conforme ABNT;
              \item Cria e alimenta a Lista de Quadros.
            \end{itemize}
        \end{itemize}
      \end{block}
    \end{column}

    \begin{column}{0.48\textwidth}
      \begin{tcolorbox}[colback=white,colframe=conepeteal,title=\textbf{Robustez Pós-Textual}]
        \begin{itemize}
          \item \textbf{Apêndices e Anexos:} Transição automática da contagem numérica para alfabética ($1, 2 \to A, B$);
          \item Não corrompe a hierarquia do sumário geral;
          \item Suporte nativo a erratas e listagens dinâmicas de equações;
          \item Compatibilização total com o pacote bibliográfico \texttt{abntex2cite}.
        \end{itemize}
      \end{tcolorbox}
    \end{column}
  \end{columns}
\end{frame}

% --- Slide 14: Conclusões e Futuro ---
\begin{frame}{Conclusões e Desdobramentos Futuros}
  \begin{block}{\textbf{Conclusões}}
    \begin{itemize}
      \item A criação da classe \texttt{ifftese.cls} e do pacote \texttt{macros.sty} cumpriu com pleno êxito a abstração das normas ABNT e IBGE;
      \item Observou-se uma redução drástica no tempo operacional de formatação e na incidência de erros de compilação entre usuários;
      \item A solução democratiza o rigor do \LaTeX\ no Instituto Federal Fluminense.
    \end{itemize}
  \end{block}

  \vspace{0.2cm}

  \begin{tcolorbox}[colback=conepeteal!10,colframe=conepeteal,title=\textbf{Trabalho em Andamento: Interface Web por Formulários}]
    \begin{itemize}
      \item Desenvolvimento de uma plataforma online simplificada (focada exclusivamente na \texttt{ifftese.cls});
      \item O autor preenche apenas formulários com metadados e textos das seções;
      \item O servidor compila e entrega o PDF diagramado com perfeição, sem necessidade de contato direto com código-fonte TeX.
    \end{itemize}
  \end{tcolorbox}
\end{frame}

% --- Slide 15: Encerramento com Fundo da Via Láctea e QR Code ---
{
\usebackgroundtemplate{%
  \includegraphics[width=\paperwidth,height=\paperheight]{vialactea_blur.jpg}%
}
\begin{frame}[plain]
  \vspace{0.6cm}
  \centering
  \begin{tcolorbox}[colback=black!75,colframe=iffgreen,width=0.88\textwidth,arc=3mm,center]
    \centering \color{white}
    {\Large \textbf{Muito Obrigado pela Atenção!}}\\[0.2cm]
    \small \textbf{CONEPE 2026} --- XIII Congresso de Ensino, Pesquisa e Extensão\\[0.3cm]
    
    \begin{columns}[c]
      \begin{column}{0.58\textwidth}
        \raggedright \small
        \textbf{Agradecimentos Institucionais:}\\
        $\bullet$ \textbf{CNPq} --- Projeto nº 129985/2025-2\\
        $\bullet$ \textbf{Instituto Federal Fluminense} --- \textit{Campus} Bom Jesus\\
        $\bullet$ \textbf{PUC-Chile} e \textbf{FAPERJ}\\[0.3cm]
        \textbf{Contato:} \texttt{pedroiff0@gmail.com}\\[0.1cm]
        \textbf{Página do projeto:} \url{https://phrandrade.com}
      \end{column}

      \begin{column}{0.35\textwidth}
        \centering
        \begin{tcolorbox}[colback=white,colframe=white,arc=2mm,width=2.6cm,center]
          \includegraphics[width=2.2cm,height=2.2cm]{qr_code_phrandrade-1.png}
        \end{tcolorbox}
        \vspace{0.1cm}
        {\scriptsize \textbf{Acesse: phrandrade.com}}
      \end{column}
    \end{columns}
  \end{tcolorbox}
\end{frame}
}

\end{document}
